\documentclass[11pt,a4paper]{article}
\usepackage[T1]{fontenc}
\usepackage[spanish]{babel}
\usepackage{csquotes}
\usepackage[margin=2cm]{geometry}
\usepackage{graphicx}
\usepackage[backend=biber,style=ieee,sorting=none]{biblatex}
\usepackage[hidelinks]{hyperref}
\addbibresource{referencias.bib}
\DefineBibliographyStrings{spanish}{artno = {art\adddot}}

\newcommand{\estudiante}{[Nombres y apellidos]}
\newcommand{\institucion}{[Institución]}
\newcommand{\curso}{[Asignatura y paralelo]}
\newcommand{\docente}{[Docente]}
\newcommand{\pendiente}[1]{\textit{[Pendiente: #1]}}

\begin{document}
\hypersetup{pageanchor=false}
\begin{titlepage}
  \centering
  {\Large \institucion\par}
  \vfill
  {\large FORM-IC-01\par}
  \medskip
  {\LARGE\bfseries Del monolito a los microservicios: el caso de Allegro\par}
  \vfill
  {\large \estudiante\par
  \curso\par
  \docente\par}
  \vfill
  {\large Noviembre de 2026\par}
  \medskip
  {\small Borrador de trabajo --- no entregar\par}
\end{titlepage}
\hypersetup{pageanchor=true}

% Reemplazar las notas de trabajo por un análisis de autoría propia,
% conservando la cita junto al dato o la afirmación que respalda.
\section*{Análisis del caso}

\paragraph{1. Caso y escala.}
Notas: plataforma de comercio electrónico de Europa Central;
20 millones de usuarios y más de 300 millones de ofertas
\cite{kosmulski2024allegro}. % Background
Más de mil servicios, según la retrospectiva de 2024
\cite{kosmulski2024allegro}. % One size does not fit all

\paragraph{2. Arquitectura anterior y motivo del cambio.}
Notas: monolito PHP, una base de datos central, conflictos de integración y
despliegues mayormente manuales. Rubicon buscó acelerar la entrega
\cite{kosmulski2024allegro}. % Background; Rubicon Rises
Paso intermedio: la plataforma SOA en PHP aún tenía un frontend compartido
y obligaba a sincronizar lanzamientos
\cite{ignasiak2015architecture}. % Switching from the old platform to the new one

\paragraph{3. Beneficios.}
Notas: ciclos de lanzamiento independientes por servicio
\cite{ignasiak2015architecture}. % Rubikon
Crear un servicio pasó de una o dos semanas a minutos; es un dato de aprovisionamiento,
no del tiempo de entrega de una funcionalidad
\cite{kosmulski2024allegro}. % Takeaways

\paragraph{4. Desafío y respuesta.}
Notas: separar los datos de la base SQL central exigió remodelar datos y operaciones
transaccionales. Ejemplo: un servicio cambió de Cassandra a MongoDB en un sprint de
dos semanas, por sus patrones de acceso y sin cambiar su API
\cite{kosmulski2024allegro}. % No-ing SQL
Contraste académico: cambios técnicos y organizativos iterativos
\cite{michaelayas2023migration}; % Sección 1
problemas de diseño y gestión de servicios observados en Stack Overflow
\cite{wu2022microservices}. % Secciones I y IV
Ejemplo posterior, si aporta al argumento: Allegro trasladó funciones de red
de las bibliotecas a una malla con Envoy mediante adopción gradual
\cite{jedrzejczyk2020mesh}. % The why; The status

\paragraph{5. Aplicación al PFC.}
\pendiente{decisión que adoptaría, qué no copiaría y por qué.}

\paragraph{Ejemplo de primera mano.}
\pendiente{identificar el sistema, cómo lo conoces y qué evidencia observaste de que sus servicios se pueden desplegar por separado.}

\paragraph{Declaración de uso de IA.}
Se utilizó IA para organizar notas con citas, contrastar datos bibliográficos
y preparar la plantilla de LaTeX. Revisar y actualizar esta declaración antes de entregar.

\begin{center}
  \IfFileExists{evidencias/allegro-anotada.png}{%
    \includegraphics[width=\linewidth,height=3.2cm,keepaspectratio]{evidencias/allegro-anotada.png}\par
    {\small Captura anotada de la fuente principal \cite{kosmulski2024allegro}.}%
  }{%
    \pendiente{incorporar la captura anotada personalmente de la fuente principal \cite{kosmulski2024allegro}.}%
  }
\end{center}

\clearpage
\printbibliography[title={Referencias}]
\end{document}
